% !TeX program = xelatex
% =============================================================
% 模板 02：中文实验报告（ctex + XeLaTeX + BibTeX / GB/T 7714）
% 编译： scripts\build.ps1 cn-lab-report
% 适用： 大学物理实验、电路实验、模电/数电实验等课程报告
% 版式： A4，上下 2.54cm，左右 3.17cm，正文小四，1.5 倍行距
% =============================================================
\documentclass[UTF8,a4paper,zihao=-4,fontset=windows]{ctexart}

\usepackage{geometry}
\geometry{top=2.54cm,bottom=2.54cm,left=3.17cm,right=3.17cm}
\usepackage{amsmath,amssymb,bm}
\usepackage{siunitx}
\usepackage{booktabs}
\usepackage{array}
\usepackage{graphicx}
\usepackage{caption}
\usepackage{gbt7714}
\usepackage[hidelinks]{hyperref}

\linespread{1.5}
\pagestyle{plain}   % 页码居中页脚，不显示页眉
\captionsetup{font=small,labelsep=quad}
\captionsetup[table]{position=top}

% 需要代码块 / 算法伪代码 / 子图 / 跨页长表 / 图表按节编号时：
% 见 references\latex-blocks.md，复制对应导言区与正文块即可

% 一级标题用“一、二、三”形式，四号黑体
\ctexset{
  section = {
    name   = {,、},
    number = \chinese{section},
    format = \zihao{4}\heiti
  },
  subsection = {format = \zihao{-4}\heiti}
}

% 实验报告常用信息栏
\newcommand{\labinfo}[4]{%
  \noindent
  % 四列宽度合计 0.86\textwidth，给 3 处列间距留出余量；p 列左对齐，避免 underfull
  \begin{tabular}{@{}>{\raggedright\arraybackslash}p{0.18\textwidth}%
                     >{\raggedright\arraybackslash}p{0.25\textwidth}%
                     >{\raggedright\arraybackslash}p{0.18\textwidth}%
                     >{\raggedright\arraybackslash}p{0.25\textwidth}@{}}
    课程名称： & #1 & 实验日期： & #2 \\
    姓\qquad 名： & #3 & 学\qquad 号： & #4 \\
  \end{tabular}\par\vspace{0.6em}\hrule\vspace{0.8em}
}

\begin{document}

% ------------------------- 标题区（改这里） -------------------------
\begin{center}
  {\zihao{3}\heiti 实验报告：共射放大电路静态工作点的测量}\\[0.6em]
  {\zihao{5}（实验报告占位标题，替换成你的实验名称）}
\end{center}

\labinfo{模拟电子技术实验}{2026-11-10}{张三}{2026000000}

% ------------------------- 一、实验目的 -------------------------
\section{实验目的}

（1）掌握共射放大电路静态工作点的测量方法；\\
（2）熟悉直流稳压电源、万用表、示波器的基本操作；\\
（3）学会由实测数据判断晶体管工作状态。

\section{实验原理}

共射放大电路的静态工作点由基极偏置电阻决定。集电极电流与集电极—发射极电压分别为
\begin{equation}
  I_{CQ}=\beta I_{BQ},\qquad
  U_{CEQ}=V_{CC}-I_{CQ}R_{C},
  \label{eq:qpoint}
\end{equation}
式中 $\beta$ 为电流放大系数，$V_{CC}$ 为电源电压，$R_{C}$ 为集电极电阻。实测中先测 $U_{BEQ}$、$U_{CEQ}$，再由式~\eqref{eq:qpoint} 反推 $I_{CQ}$。

\section{实验器材}

直流稳压电源 1 台，数字万用表 1 台，示波器 1 台，面包板 1 块，晶体管 9013 一只，电阻、电容若干。

\section{实验内容与步骤}

（1）按原理图搭接电路，检查无误后接通电源；\\
（2）调节 $R_{b}$，使 $U_{CEQ}\approx V_{CC}/2$，记录此时各点电压；\\
（3）改变 $R_{b}$，重复测量 3 组数据；\\
（4）用示波器观察输出波形，判断是否出现饱和或截止失真。

\section{数据记录与处理}

测量数据见表~\ref{tab:data}，表中 $U_{BEQ}$、$U_{CEQ}$ 由万用表直流电压挡测得，$I_{CQ}$ 由式~\eqref{eq:qpoint} 计算。

\begin{table}[htbp]
  \centering
  \caption{静态工作点测量数据}
  \label{tab:data}
  \begin{tabular}{cccccc}
    \toprule
    序号 & $R_{b}$/k$\Omega$ & $U_{BEQ}$/V & $U_{CEQ}$/V & $I_{CQ}$/mA & 工作状态 \\
    \midrule
    1 & 300 & 0.68 & 5.98 & 1.01 & 放大区 \\
    2 & 200 & 0.70 & 3.12 & 1.94 & 放大区 \\
    3 & 100 & 0.72 & 0.35 & 2.47 & 饱和区 \\
    \bottomrule
  \end{tabular}
\end{table}

以第 2 组数据为例，相对误差按
\begin{equation}
  \delta=\frac{\left|I_{CQ,\text{实}}-I_{CQ,\text{理}}\right|}{I_{CQ,\text{理}}}\times100\%
  \label{eq:error}
\end{equation}
计算，得 $\delta\approx 3.2\%$，在允许范围内。

\section{误差分析}

主要误差来源：万用表内阻分流、电阻标称值与实际值偏差、晶体管 $\beta$ 随温度漂移。改进办法是采用高输入阻抗仪表并在恒温环境下测量。

\section{结论}

通过本次实验测得三组静态工作点数据，验证了式~\eqref{eq:qpoint} 的正确性；当 $R_{b}$ 减小到 100\,k$\Omega$ 时晶体管进入饱和区，与理论分析一致\cite{zhang2024demo}。

\bibliographystyle{gbt7714-numerical}
\bibliography{refs}

\end{document}
